\documentclass[10pt,a4paper]{amsart}
\usepackage[margin=19mm]{geometry}
\usepackage{amsmath,amssymb,mathtools}
\usepackage[scr=rsfs,cal=euler]{mathalfa}
\usepackage{xfrac}
\usepackage[unicode,hidelinks,bookmarks=false]{hyperref}
\newcommand{\ie}{i.e.\ }
\newcommand{\cf}{cf.\ }
\newcommand{\resp}{respectively\ }
\newcommand{\Subsection}{\subsection}
\providecommand{\nobreakdash}{\nobreak\hbox{-}}
\setlength{\emergencystretch}{2em}
\multlinegap0pt
\usepackage{bm}
\usepackage{bbm}
\usepackage{dsfont}
\usepackage{xspace}
\usepackage{cancel}
\usepackage{mathtools}
\usepackage{amsthm}
\makeatletter
\@ifundefined{theorem}{\newtheorem{theorem}{Theorem}}{}
\@ifundefined{problem}{\newtheorem{problem}[theorem]{Problem}}{}
\makeatother
\newcommand{\R}{\mathbb{R}}
\newcommand{\Z}{\mathbb{Z}}
\newcommand{\mc}[1]{\mathcal{#1}}
\newcommand{\rs}{\setminus}
\title[Eulerian-spanning set and coboundary operator: An investigat]{Eulerian-spanning set and coboundary operator: An investigation of maxcut beyond planar graphs}
\author{Qiming Fang, Sihong Shao, Yuxuan Wu}
\date{Source: arXiv:2606.00725v1 (CC BY 4.0)}
\begin{document}
\maketitle

\noindent\emph{Source and licence.} Eulerian-spanning set and coboundary operator: An investigation of maxcut beyond planar graphs. Version arXiv:2606.00725v1 (CC BY 4.0), distributed under the Creative Commons Attribution 4.0 International licence, as recorded in the arXiv licence metadata for this exact version and shown on the abstract page https://arxiv.org/abs/2606.00725v1. Full rights notice: licenses/arxiv-2606.00725-CC-BY-4.0.txt.\par\smallskip

\noindent\textit{Editorial scope.} Counted unit: the Theorem whose source label is \texttt{None} in the article (source0.tex lines 569--571, source label \texttt{None}), delivered together with the article's own complete original proof of it (source0.tex lines 572--603, one \texttt{proof} environment of 32 lines). No mathematics is written, repaired or re-derived: statement and proof are the article's own text line for line. Required context retained uncounted: prb31 (lines 237--239); prb32 (lines 241--243); thm41 (lines 558--560). Every definition, display and label of those lines is retained unchanged. Omitted from this package and not redistributed: 1--236, 240--240, 244--557, 561--568, 604--658. Nothing inside a retained excerpt is omitted or altered: every retained body is delivered exactly as the article's lines are written, so no omission receipt of any kind accompanies this package. Citations: the retained excerpts carry no citation command at all, so no bibliography entry of the article is transcribed and no reference is redistributed. Reference adaptation: none was needed and none was made; every reference inside the delivered unit resolves inside it, so no unresolved reference or citation is produced. Printed slips kept literal: no printed word, symbol, hypothesis or number of the counted statement or of its proof was altered. Layout and class: the article's own class is \texttt{article} with packages including amsmath, amssymb, amsfonts, appendix, amsthm, extarrows, enumerate, graphicx, eepic, color, colordvi, amscd, subcaption, tikz; the wrapper is the lane's pinned \texttt{amsart} editorial wrapper at 10pt on A4, which restates the theorem-like environments the retained text needs and adds the article's own macro definitions that the retained text needs (\texttt{\textbackslash R}, \texttt{\textbackslash Z}, \texttt{\textbackslash mc}, \texttt{\textbackslash rs}), with extra wrapper packages (bm, bbm, dsfont, xspace, cancel, mathtools). The delivered numbering is the unit's own. Assets: no retained span contains a figure environment, a graphics-inclusion command, a PDF-page inclusion, a table or a TikZ picture; no image file is copied into this package, the package declares no asset dependency and no image file is redistributed with it. Licence: version arXiv:2606.00725v1 of this article is distributed under the Creative Commons Attribution 4.0 International licence, as recorded in the arXiv licence metadata for this exact version and shown on the abstract page https://arxiv.org/abs/2606.00725v1. A full rights notice is delivered at licenses/arxiv-2606.00725-CC-BY-4.0.txt. Characters: every retained excerpt is pure ASCII, so the delivered engine, XeLaTeX with its default Latin Modern text font, reports no missing character.
	\begin{problem}[\it{WSSD}]\label{prb31}
		Suppose $S$ is a finite multiset, and $\mathcal{S}=\{S_i\}_{i=1}^{m}\subseteq\Z_2^S$ is a multiset consisting of subsets of $S$. Moreover, there is a weight function $w:\mc{S}\to\R^+\cup\{0\}$. Suppose we are given a target set $T\subseteq S$. The goal is to find a multiset $\mathcal{T}\subseteq\mc{S}$, such that $\triangle_{S_i\in\mc{T}}S_i=T$, and the total weight of sets in $\mc{T}$ is minimal, or claim that no such $\mc{T}$ exists.
	\end{problem} \par

	\begin{problem}[\it{decision version of Problem \ref{prb31}}]\label{prb32}
		Suppose $S$ is a finite multiset, and $\mathcal{S}=\{S_i\}_{i=1}^{m}\subseteq\Z_2^S$ is a multiset consisting of subsets of $S$. Moreover, there is weight function $w:\mc{S}\to\R^+\cup\{0\}$. Suppose we are given a target set $T\subseteq S$, and a threshold real number $r$. The goal is to decide, whether there is a multiset $\mathcal{T}\subseteq\mc{S}$, such that $\triangle_{S_i\in\mc{T}}S_i=T$, and the total weight of sets in $\mc{T}$ is no more than $r$.
	\end{problem}

	\begin{theorem}\label{thm41}
		Problem \ref{prb32} under the restriction that $w\equiv 1$ is NP-complete.
	\end{theorem}

	\begin{theorem}
		Problem \ref{prb32} under the restriction that $w\equiv1$, and $\forall S_i\in\mc{S},|S_i|\in\{2,3\}$ is NP-complete. 
	\end{theorem}

	\begin{proof}
		To see that the problem is NP-hard, we reduce the problem in Theorem \ref{thm41} to this problem. We only need to take care of sets of size 1 or at least 4. We break the reduction into three steps: \par
		Step 1. In step 1, we eliminate all subsets of size 1. Suppose $S,\mc{S},T,r$ is an instance of the problem, and there is a set $S_i\in\mc{S}$ of size 1, and $S_i=\{s\}$, we do some modifications to the problem. Let
		$$S'=S\rs\{s\}\cup\{s^{(1)},s^{(2)}\}$$ 
		$$\mc{S}'=\{S_1',\cdots,S_m\},\ 
		S_k'=\left\{
		\begin{aligned}
			&S_k,\ if\ s\notin S_k \\
			&S_k\rs\{s\}\cup\{s^{(1)},s^{(2)}\}\ if\ s\in S_k
		\end{aligned}
		\right.
		$$
		$$T'=\left\{
		\begin{aligned}
			&T,\ if\ s\notin T \\
			&T\rs\{s\}\cup\{s^{(1)},s^{(2)}\}\ if\ s\in T
		\end{aligned}
		\right.
		$$ \par
		Then $S',\mc{S}',T',r$ is another instance of the problem in Theorem \ref{thm41}, which is equivalent to $S,\mc{S},T,r$. Moreover, the number of sets of size 1 in $\mc{S}'$ is strictly smaller than that of $\mc{S}$. Therefore, by repeating the above procedure, we will finally get an instance $S^{(1)},\mc{S}^{(1)},T^{(1)},r$ such that all sets in $\mc{S}^{(1)}$ are of size at least 2. \par
		Step 2. We give every set in $\mc{S}^{(1)}$ a weight of $p=|S^{(1)}|+2$, and denote the weight function by $w$. Now we have an instance, i.e. $S^{(1)},\mc{S}^{(1)},T^{(1)},w,pr$, which is equivalent to $S^{(1)},\mc{S}^{(1)},T^{(1)},r$. By a slight abuse of notation, we will denote this instance by $S,\mc{S},T,w,pr$, where $w\equiv|S|+2,p=|S|+2$.\par 
		Step 3. Decompose all subsets of size at least 4. We actually decompose every set in $\mc{S}$ in this step. Suppose $S_i\in\mc{S},S_i=\{s_i\}_{i=1}^h,w(S_i)=p$. Because $h\leq p-2$, we can do the following modifications: let
		$$S'=S\cup\{z_1,z_2,\cdots,z_{p-1},z_p\}$$
		$$\mc{S}'=\mc{S}\rs\{S_i\}\cup\mc{A}$$
		$$\mc{A}=\mc{A}_1\cup\mc{A}_2\cup\mc{A}_3$$
		$$\mc{A}_1=\{\{s_1,z_1\},\{s_2,z_1,z_2\},\{s_3,z_2,z_3\},\cdots,\{s_h,z_{h-1},z_h\}\}$$
		$$\mc{A}_2=\{\{z_h,z_{h+1}\},\cdots,\{z_{p-3},z_{p-2}\}\},\quad (\text{if }h=p-2,\mc{A}_2=\emptyset)$$
		$$\mc{A}_3=\{\{z_{p-2},z_{p-1},z_p\},\{z_{p-1},z_p\}\}$$
		$$\forall j\neq i, w'(S_j)=w'(S_i)$$
		$$\forall A\in\mc{A},w(A)=1$$
		Then $S',\mc{S}',T,w',pr$ is another instance equivalent to $S,\mc{S},T,w,pr$ . By repeating the above procedure, for every $S_i\in\mc{S},S_i=\{s_i\}_{i=1}^h,w(S_i)=p$, it is decomposed to sets of size 2 or 3, and of weight 1. Therefore we finally get an instance $S^{(2)},\mc{S}^{(2)},T,pr$ such that every set in $\mc{S}^{(2)}$ is of size 2 or 3. The reduction is complete, and can be performed in time polynomial to the size of $|S|$ before the reduction.
	\end{proof}

\end{document}
